%-------------------------
% Resume in Latex
% Author : Chandula Adhikari
% Based on: https://github.com/jakegut/resume (which was itself based on https://github.com/sb2nov/resume)
% License : MIT
%------------------------

\documentclass[a4paper,10pt]{article}
\usepackage{latexsym}
\usepackage[empty]{fullpage}
\usepackage{titlesec}
\usepackage{marvosym}
\usepackage[usenames,dvipsnames]{color}
\usepackage{verbatim}
\usepackage{enumitem}
\usepackage[hidelinks]{hyperref}
\usepackage{fancyhdr}
\usepackage[english]{babel}
\usepackage{multicol}
\usepackage{tabularx}
\usepackage[dvipsnames]{xcolor}
\usepackage{array}
\usepackage{tabularx}
\usepackage{fontawesome5}
\setlength{\multicolsep}{-3.0pt}
\setlength{\columnsep}{-1pt}
\setlength{\topmargin}{0.55in}

\input{glyphtounicode}

% -------------------- FONT OPTIONS --------------------
% sans-serif
% \usepackage[sfdefault]{roboto}
% \usepackage[sfdefault]{noto-sans}
% serif
% \usepackage{charter}

\pagestyle{fancy}
\fancyhf{} % clear all header and footer fields
\fancyfoot{}
\renewcommand{\headrulewidth}{0pt}
\renewcommand{\footrulewidth}{0pt}

% Adjust margins
\addtolength{\oddsidemargin}{-0.5in}
\addtolength{\evensidemargin}{-0.5in}
\addtolength{\textwidth}{1in}
\addtolength{\topmargin}{-1in} % Default was -.5in
\addtolength{\textheight}{1.0in}

\urlstyle{same}

\raggedbottom
\raggedright
\setlength{\tabcolsep}{0in}

% Section formatting
\titleformat{\section}{
  \vspace{-5pt}\color{MidnightBlue}\scshape\raggedright\large
}{}{0em}{}[\color{gray}\titlerule \vspace{-5pt}]



% Subsection formatting
\titleformat{\subsection}{
  \vspace{-4pt}\scshape\raggedright\large
}{\hspace{-.15in}}{0em}{}[\color{gray}\vspace{-8pt}]

% Ensure that generate pdf is machine readable/ATS parsable
\pdfgentounicode=1

% -------------------- CUSTOM COMMANDS --------------------
\newcommand{\resumeItem}[1]{
  \item\small{
    {#1 \vspace{-2pt}}
  }
}

\newcommand{\resumeSubheading}[4]{
  \vspace{-2pt}\item
    \begin{tabular*}{0.97\textwidth}[t]{l@{\extracolsep{\fill}}r}
      \textbf{#1} & \footnotesize{#2} \\
      \small#3 & \small #4 \\
    \end{tabular*}\vspace{-7pt}
}

\newcommand{\resumeSubheadingSub}[6]{
  \vspace{-2pt}\item
    \begin{tabular*}{0.97\textwidth}[t]{l@{\extracolsep{\fill}}r}
      \textbf{#1} & \footnotesize{#2} \\
      {\small#3} & {\small #4} \\
      \textit{\small#5} & \textit{\small #6} \\
    \end{tabular*}\vspace{-7pt}
}

\newcommand{\resumeSubSubheading}[2]{
    \item
    \begin{tabular*}{0.97\textwidth}{l@{\extracolsep{\fill}}r}
      \textit{\small#1} & \textit{\small #2} \\
    \end{tabular*}\vspace{-7pt}
}

\newcommand{\resumeSkillSubheading}[2]{
    \vspace{-2pt}\item
    \begin{tabularx}{\textwidth}{@{}>{\hsize=.40\hsize}X>{\hsize=1.60\hsize}X@{}}
      \textbf{\small#1} & \small #2 \\
    \end{tabularx}\vspace{-7pt}
}

\newcommand{\resumeExperienceSubheading}[2]{
    \vspace{-19pt}\item
    \begin{tabularx}{\textwidth}{@{}>{\hsize=1.4\hsize}X>{\hsize=0.6\hsize}X@{}}
      \small#1 & \small #2 \\
    \end{tabularx}\vspace{-7pt}
}


\newcommand{\resumeProjectHeading}[2]{
    \item
    \begin{tabular*}{0.97\textwidth}{l@{\extracolsep{\fill}}r}
      #1 & \footnotesize #2 \\
    \end{tabular*}\vspace{-7pt}
}

\newcolumntype{B}{>{\bfseries}l}
\newcolumntype{L}{X}

\newcommand{\resumeSubItem}[1]{\resumeItem{#1}\vspace{-4pt}}
\newcommand{\resumeSubHeadingListStart}{\begin{itemize}[leftmargin=0.15in, label={}]}
\newcommand{\resumeSubHeadingListEnd}{\end{itemize}}
\newcommand{\resumeItemListStart}{\begin{itemize}}
\newcommand{\resumeItemListEnd}{\end{itemize}\vspace{-5pt}}

\renewcommand\labelitemii{$\vcenter{\hbox{\tiny$\bullet$}}$}

\setlength{\footskip}{4.08003pt}



% -------------------- START OF DOCUMENT --------------------
\begin{document}

% -------------------- HEADING--------------------
\begin{flushright}
  %\vspace{-4pt}
  \color{gray}
  \item
  % Last Updated on March 28th, 2024
\end{flushright}

\vspace{-0.6in}
\begin{center}
    \textbf{\Huge \scshape Chandula Adhikari} \\ \vspace{2pt}
    \large
    \textcolor{MidnightBlue}{Department of Computer Engineering, University of Peradeniya, Sri Lanka}\\ \vspace{5pt}
    \small
    \faIcon{envelope}
    \href{mailto:rajcadhikari@gmail.com}
    {\underline{rajcadhikari@gmail.com}} $  $ $ \mid $ $  $
    \faIcon{phone}
    +94 711530046 $  $ $ \mid $ $  $
    \faIcon{github}
    \href{https://github.com/chandula00}{\underline{github.com/chandula00}} \\ \vspace{2pt}
    \faIcon{linkedin}
    \href{https://www.linkedin.com/in/janith-chandula-adhikari-5543a5223}{\underline{linkedin.com/in/chandula-adhikari}} $  $ $ \mid $ $  $
    \faIcon{globe}
    \href{https://chandula00.github.io/}{\underline{chandula00.github.io}}
\end{center}

% -------------------- ABOUT ME --------------------
\section{\textbf{Research Interests}}
    \vspace{1pt}
    % \begin{center}
    %     \hspace{0.12cm}
    %     Computer Vision \hspace{2.13cm}
    %     Machine Learning \hspace{2.13cm}
    %     Explainable AI \hspace{2.13cm}
    %     Embedded Systems
    % \end{center}



        I work on machine learning \textbf{shaped by the biology it describes}, rather than generic architectures borrowed from other domains. My focus is \textbf{single-cell genomics}: how genes are regulated, how cells change state, and what can be recovered from sparse, noisy measurements. More broadly, I am drawn to \textbf{generative modelling across biomedical data}. I want a model's assumptions to stay clear enough for biologists and clinicians to question them.







% -------------------- EDUCATION --------------------
\section{\textbf{Education}}
    \resumeSubHeadingListStart

        \vspace{2pt}
        \resumeSubheadingSub
            {University of Peradeniya}{\textcolor{MidnightBlue}{Mar. 2021 - Aug. 2025}}
            {BSc.Eng(Hons.) in Computer Engineering --- \textbf{First Class}}{Overall GPA: \textbf{\normalsize 3.85/4.0}}
            {}{\textup{Final-Year GPA: \textbf{\normalsize 4.00/4.00}}}
        \vspace{2pt}
        \resumeSubheadingSub
            {Dharmaraja College, Kandy}{\textcolor{MidnightBlue}{Jan. 2006 - Dec. 2019}}
            {G.C.E. Advanced Level Examination - 467 from 19500+ participants nationwide}{Z-Score: \textbf{\normalsize 2.2069}}{}{}
        % {}{467 from 35000+ participants nationwide}

        % \subsection{Coursework}
        %   \textbf{Courses:} Object-Oriented Programming, Data Structures \& Algorithms, Embedded Systems, Discrete Math, Linear Algebra, Calculus, Physics, Probability \& Statistics \\
        %   \textbf{Awards:} Dean's Honor List (3x), AP Scholar with Distinction (2x), World Language History Award (Spanish)

    \resumeSubHeadingListEnd
    \vspace{-20pt}

% % -------------------- Publications --------------------
\section{\textbf{Publications}}
  \resumeSubHeadingListStart
        \vspace{2pt}

        \resumeProjectHeading
            {\textbf{\href{https://chandula00.github.io/TARVI/}
            {\underline{TARVI: Transcription-Factor Aided RNA Velocity Inference with Supervised Latent Time}}}}{}
            \vspace{-2pt}
            \resumeExperienceSubheading
                {\emph{\textbf{R.A.J.C. Adhikari}, S. Dassanayake, D. Herath}}{}
            \vspace{-12pt}
            \resumeItemListStart
            \resumeItem{Replaced the static per-gene transcription rate used by existing RNA velocity models with \textbf{cell-specific, TF-regulated rates}, injecting a masked learnable TF$\rightarrow$gene matrix derived from \textbf{ENCODE ChIP-seq and ChEA} regulatory priors into a variational autoencoder --- a knowledge-guided constraint that makes the model's regulatory assumptions explicit and inspectable.}
            \resumeItem{Introduced a supervised latent-time head trained with ODE-residual, self-distillation and graph-Laplacian smoothness losses, plus post-hoc velocity--pseudotime blending that grounds local velocity vectors in global trajectory topology.}
            \resumeItem{Improved pseudotime calibration by \textbf{126\%} over VeloVI (Spearman 0.747 vs. 0.331) across 5 single-cell datasets (pancreas, bone marrow, forebrain, chromaffin, scEU organoid), while matching the strongest of 4 baselines on directional metrics.}
            \resumeItem{\textbf{Best Paper Award} --- Biomedical Engineering \& Instrumentation track $|$ MERCon 2026 (IEEE) $|$ Code: \href{https://github.com/chandula00/TARVI}{\underline{github.com/chandula00/TARVI}}}
          \resumeItemListEnd

        \vspace{4pt}
        \resumeProjectHeading
            {\textbf{Dropout Is the Absorbing State: Discrete Diffusion for Single-Cell Gene}}{}
            \vspace{-14pt}
        \resumeProjectHeading
            {\textbf{Expression Imputation} $|$ \emph{\footnotesize Under review, AAAI 2027}}{}
            \vspace{-2pt}
            \resumeExperienceSubheading
                {\emph{\textbf{R.A.J.C. Adhikari} et al.}}{}
            \vspace{-12pt}
            \resumeItemListStart
            \resumeItem{Reformulated scRNA-seq dropout as the \textbf{absorbing state of a discrete diffusion process}: expression is tokenized into per-gene quantile bins and only candidate dropout entries are corrupted toward a reserved \texttt{[MASK]} state, so measured entries stay immutable, non-negativity holds by construction, and each gene receives a full categorical distribution rather than a conditional mean.}
            \resumeItem{Designed a role-conditioned Transformer denoiser and a \textbf{confidence-first parallel decoding} scheme that exploits the monotonicity of the absorbing chain to reconstruct the matrix in a few passes instead of a full ancestral trajectory.}
            \resumeItem{Against \textbf{9 baselines} (MAGIC, scImpute, SAVER, ALRA, DCA, scVI, scIGANs, scIDPMs) on 4 simulated datasets at 50--72\% dropout, achieved the best rank correlation ($\rho$ 0.90--0.94) and lowest log-scale RMSE (0.55--0.61) on every dataset, and won \textbf{20 of 21} cell-type clustering views across 3 real datasets.}
          \resumeItemListEnd

        \vspace{4pt}
        \resumeProjectHeading
            {\textbf{\href{https://doi.org/10.1109/ICATC68823.2025.11407838}{\underline{Uncovering Hidden Temporal Patterns in T20I Cricket Through Ball-by-Ball Data Analysis}}}}{}
            \vspace{-2pt}
            \resumeExperienceSubheading
                {\emph{\textbf{R.A.J.C. Adhikari} et al.}}{}
            \vspace{-12pt}
            \resumeItemListStart
            \resumeItem{Analyzed 165,000+ deliveries from 1,200+ T20I matches (2006-2025) to identify temporal patterns in batting and bowling performance across powerplay, middle, and death overs.}
            \resumeItem{ICATC 2025 $|$ IEEE Xplore}
          \resumeItemListEnd

        \vspace{4pt}
        \resumeProjectHeading
            {\textbf{\href{https://doi.org/10.1109/ICAC64487.2024.10851123}
            {\underline{Optimized Multi-Processor System-on-Chip (MPSoC) Design for Low-Resource JPEG Encoding}}}}{}
            \vspace{-2pt}
            \resumeExperienceSubheading
                {\emph{K.H. Gunawardana, \textbf{R.A.J.C. Adhikari}, I. Nawinne}}{}
            \vspace{-12pt}
            \resumeItemListStart
            \resumeItem{Proposed a pipelined MPSoC architecture for efficient JPEG encoding using Altera Nios II/e cores on a Cyclone IV FPGA, integrating custom instructions, FIFO-based inter-core communication, and superscalar enhancements to maximize throughput.}
            \resumeItem{ICAC 2024 $|$ IEEE Xplore}
          \resumeItemListEnd
    \resumeSubHeadingListEnd
    \vspace{-10pt}

% -------------------- EXPERIENCE --------------------
\section{\textbf{Experience}}
  \resumeSubHeadingListStart
        \vspace{2pt}

        \resumeProjectHeading
            {\textbf{Research Assistant} $|$
            \footnotesize\emph{\href{https://www.linkedin.com/company/multidisciplinaryai}{\underline{Multidisciplinary AI Research Centre (MARC)}}, University of Peradeniya} \vspace{5pt}}{\textcolor{MidnightBlue}{June 2026 -- Present}}

                \small{Research post at the University of Peradeniya's interdisciplinary AI centre, which coordinates machine learning across engineering, medicine and the life sciences.\\
                \textbf{Lead two research groups} at the centre: the \textbf{Computational Biology \& Bioinformatics} pillar, and \textbf{Generative AI for Dermatology}.\\
                Work carried out in this post produced the absorbing-state discrete diffusion imputation model currently \textbf{under review at AAAI 2027}.}{}

        \vspace{4pt}
        \resumeProjectHeading
            {\textbf{\href{https://people.ce.pdn.ac.lk/staff/temporary-academic-staff/}
            {\underline{Temporary Instructor}}} $|$
            \footnotesize\emph{Department of Computer Engineering} \vspace{5pt}}{\textcolor{MidnightBlue}{Sep. 2025 -- June 2026}}

                \small{Conducted laboratory sessions and prepared tutorials and lab specifications for undergraduate courses; supervised and mentored student project work. Built on three prior years as an Undergraduate Teaching Assistant (Nov. 2023 -- July 2025) across Data Structures \& Algorithms, Computer Architecture, Digital Design and Embedded Systems.}{}

        \vspace{4pt}
        \resumeProjectHeading
            {\textbf{Intern Data Scientist} $|$
            \footnotesize\emph{\href{https://www.octave.lk/}{\underline{OCTAVE}}, John Keells Holdings} \vspace{5pt}}{\textcolor{MidnightBlue}{July 2024 -- Dec. 2024}}

                \small{Built and validated production-grade data pipelines over large enterprise datasets, and developed customer segmentation models that directly supported live digital marketing campaigns --- experience in end-to-end data science workflow, reproducibility and experiment tracking that transfers to large-scale genomic data analysis.\\}{}
            \small
                {\textbf{Technologies:} Databricks, PySpark, MLflow, Azure DevOps}{}
            \vspace{-4pt}

    \resumeSubHeadingListEnd

% -------------------- PROJECTS --------------------
\section{\textbf{Selected Projects}}
    \resumeSubHeadingListStart
        \vspace{4pt}
        \resumeProjectHeading
        {\textbf{AI-Assisted Tool for Assessing Quality of Final Root Canal Treatment Using}}{\textcolor{MidnightBlue}{Jan. 2025 -- Present}}
        \resumeProjectHeading
            {\textbf{Dental IOPA Radiographs} $|$ \emph{Group $|$ \href{https://github.com/cepdnaclk/e19-4yp-Quality-Assessment-of-Final-Root-Canal-Treatment-Using-Intra-Oral-Periapical-Radiographs}{\faIcon{github}}}}{}
        \resumeItemListStart
            \resumeItem{Developing an explainable AI framework for automated root canal treatment (RCT) quality assessment using 1,000 anonymized intraoral periapical radiographs. The four-stage pipeline—comprising dataset curation, expert-guided annotation, segmentation, and multi-parameter evaluation—enables pixel- and geometry-based analysis of clinical parameters such as filling length, lateral seal, voids, and irregularities.}
            \resumeItem{\textbf{Contribution:} Dataset preparation, annotation design, quality parameter evaluation, and explainable output visualization.}
            \resumeItem{\textbf{Technologies:} Python, PyTorch, YOLO, OpenCV, Keras.}
        \resumeItemListEnd

        \vspace{4pt}
        \resumeProjectHeading
        {\textbf{Reliability of RNA Velocity in Low-Dimensional Projections} $|$ \emph{Individual}}{\textcolor{MidnightBlue}{Feb. 2025 -- Dec. 2025}}
        \resumeItemListStart
            \resumeItem{Benchmarked custom autoencoder architectures against standard embeddings (PCA, t-SNE, UMAP) across 4 biological datasets, showing that denoising autoencoders better preserve local neighborhood structure and velocity-field coherence --- establishing that the embedding, not the velocity estimator alone, limits trajectory inference.}
            \resumeItem{Findings motivated and fed directly into \textbf{TARVI} (MERCon 2026, Best Paper Award), which addressed the resulting calibration gap at the model level.}
            \resumeItem{\textbf{Technologies:} Denoising Autoencoders, VAEs $|$ PCA, t-SNE, UMAP $|$ scVelo (dynamical model), Velocyto $|$ Evaluation: k-NN preservation, distance correlation, Jaccard similarity}
        \resumeItemListEnd


        \vspace{4pt}
        \resumeProjectHeading
            {\textbf{BeeZee: Smart Beehive Monitoring System} $|$ \emph{Group $|$ \href{https://github.com/cepdnaclk/e19-3yp-beehive-monitoring-system}{\faIcon{github}} $|$ \href{https://cepdnaclk.github.io/e19-3yp-beehive-monitoring-system/}{\faIcon{globe}}}}{\textcolor{MidnightBlue}{Nov. 2023 - July 2024}}
        \resumeItemListStart
            \resumeItem{Collaborative research with the Faculty of Agriculture, University of Peradeniya to detect early signs of bee colony abscondment. Designed and built a cloud-connected data collection unit with sensors ($CO_2$, humidity, temperature, weight) and camera, integrated into a real-time dashboard. Successfully implemented pollen-carrying bee detection using YOLOv8 and StrongSORT for object tracking.}
            \resumeItem{\textbf{Contribution:} Hardware enclosure, circuit assembly, sensor calibration, the front-end dashboard, and the computer vision pipeline for bee detection and tracking.}
            \resumeItem{\textbf{Technologies:} Raspberry Pi, AWS(S3, IoT Core, Lambda), Node.js, MongoDB, YOLOv8, StrongSORT.}
          \resumeItemListEnd


        % \vspace{4pt}
        % \resumeProjectHeading
        %     {\textbf{CricVision: Optimizing Dynamic Batting Orders in T20 Cricket} $|$ \emph{Group $|$ \href{https://github.com/cepdnaclk/e19-co544-cricket-analytics-and-prediction}{\faIcon{github}} $|$ \href{https://cepdnaclk.github.io/e19-co544-cricket-analytics-and-prediction/}{\faIcon{globe}}}}{\textcolor{MidnightBlue}{Mar. 2024 -- Present}}
        % \resumeItemListStart
        %     \resumeItem{Developed a data-driven system to dynamically optimize batting orders in T20 cricket using machine learning. Built a multi-output regression pipeline to predict batter performance (runs, strike rate), final team score, and Net Run Rate (NRR), enhancing decision-making after each wicket.}
        %     \resumeItem{\textbf{Data set :} \href{https://cricsheet.org/}
        %     {\underline{cricsheet.org}}}
        %     \resumeItem{\textbf{Contribution:} Engineered features from Cricsheet data, web-scraped contextual data (player types, weather), implemented multiple regression models (XGBoost, Chained, Neural Networks), and deployed the best-performing model via Flask for real-time insights.}
        %     \resumeItem{\textbf{Technologies:} Python, Scikit-learn, XGBoost, Flask, Pandas, Matplotlib, BeautifulSoup.}
        % \resumeItemListEnd


        % \vspace{4pt}
        % \resumeProjectHeading
        %     {\textbf{VisitLog: Digital Reporting Platform for Technical Visits} $|$ \emph{Group $|$ \href{https://github.com/cepdnaclk/e19-co227-digital-reporting-of-technical-visits}{\faIcon{github}} $|$ \href{https://cepdnaclk.github.io/e19-co227-digital-reporting-of-technical-visits/}{\faIcon{globe}}}}{\textcolor{MidnightBlue}{Aug. 2023 -- Nov. 2023}}
        % \resumeItemListStart
        %     \resumeItem{Developed a cross-platform system to streamline technical service documentation, covering repairs, troubleshooting, and maintenance, with recipient approval. The solution includes a Flutter-based mobile app for technicians and a React.js dashboard for administrators to manage jobs and personnel.}
        %      \resumeItem{\textbf{Contribution:} Built the entire mobile app, designed and developed the admin dashboard, managed Firebase collections, and implemented an NLP-based technician allocation model using availability, expertise, and job descriptions.}
        %     \resumeItem{\textbf{Technologies:} React.js, Flutter, Firebase, Figma, Python.}
        % \resumeItemListEnd

        % \vspace{4pt}
        % \resumeProjectHeading
        %     {\textbf{Obstacle Robot Swarm for Swarm Robotic Project} $|$ \emph{Group $|$ \href{https://github.com/Pera-Swarm}{\faIcon{github}} $|$ \href{https://pera-swarm.ce.pdn.ac.lk/}{\faIcon{globe}}}}{\textcolor{MidnightBlue}{Mar. 2024 -- Present}}
        % \resumeItemListStart
        %     \resumeItem{An obstacle bot is an automated robot designed for a swarm robotic arena. It can move to desired positions, avoid collisions, and be programmed as static or dynamic obstacles.}
        %     \resumeItem{ \textbf{Contribution:} Updated firmware to support autonomous collision handling and integrated obstacle robots into the existing swarm platform for simulating dynamic obstacle scenarios.}
        %     \resumeItem{\textbf{Technologies:} Arduino, Python, Java}
        % \resumeItemListEnd


        % \vspace{8pt}
        % \resumeProjectHeading
        %     {\textbf{IntelliSwitcher: Intelligent Domestic Energy Saving System} $|$ \emph{Group $|$ \href{https://github.com/IntelliSwitcher}{\faIcon{github}}}}{\textcolor{MidnightBlue}{June 2023 -- Sep. 2023}}
        % \resumeItemListStart
        %     \resumeItem{IoT device and mobile app system designed to empower households to effortlessly reduce energy consumption}
        %     \resumeItem{By continuously monitoring power factors, voltage, and current consumption, we generate real-time predictions of monthly energy consumption.}
        %     \resumeItem{\textbf{Technologies:} React Native, Firebase, scikit-learn}
        % \resumeItemListEnd

        % \vspace{4pt}
        % \resumeProjectHeading
        %     {\textbf{MPSoC Design for JPEG encoding} $|$ \emph{Group $|$ \href{https://github.com/chandula00/CO503-Advanced-Embedded-Systems-Labs}{\faIcon{github}}}}{\textcolor{MidnightBlue}{May. 2024 -- Present}}
        % \resumeItemListStart
        %     \resumeItem{Implementing a pipelined JPEG encoder on an FPGA(Altera DE2-115) using Nios II CPUs and hardware FIFO queues. The JPEG encoding process is broken down into six stages}
        %     \resumeItem{Using custom instruction hardware blocks, Custom FIFOs and implementing a super-scalar pipeline improved the performance}
        %     \resumeItem{\textbf{Technologies:} Quartus II , Nios II Software Build Tool, Altera DE2-115}
        % \resumeItemListEnd


        \vspace{2pt}

    \resumeSubHeadingListEnd




% -------------------- SKILLS --------------------
\section{\textbf{Technical Skills}}
\vspace{4pt}
\begin{itemize}[leftmargin=0.15in, label={}]
\item
\begin{tabularx}{0.97\textwidth}{@{}>{\bfseries\small\raggedright\arraybackslash}p{1.45in}>{\small}X@{}}
  Languages & Python, C/C++, Java, JavaScript/TypeScript, Verilog HDL, ARM Assembly \\[3pt]
  Machine \& Deep Learning & PyTorch, TensorFlow/Keras, scikit-learn, Hugging Face Transformers \\[3pt]
  Model Classes & Transformers, Diffusion Models (DDPM/DDIM, absorbing-state discrete diffusion), VAEs \& Autoencoders, CNNs, ODE-constrained neural models \\[3pt]
  Single-Cell \& Bioinformatics & Scanpy, AnnData, scVelo, Velocyto, Seurat; ENCODE ChIP-seq \& ChEA regulatory atlases; trajectory inference, dropout imputation, cell-type clustering \\[3pt]
  Data \& Cloud & NumPy, Pandas, Matplotlib, Seaborn, OpenCV; Databricks, PySpark, SQL (Spark SQL \& T-SQL), MLflow, Azure DevOps, AWS, Power BI \\[3pt]
  Developer Tools & Git/GitHub, Unix Shell, Linux/HPC, VS Code, PyCharm/IntelliJ, \LaTeX, Jekyll, Jenkins \\[3pt]
  Web Frameworks & React.js, Express.js, Flutter, Spring-Boot, FastAPI \\
\end{tabularx}
\end{itemize}
\vspace{-8pt}

% -------------------- ACHIEVEMENTS --------------------
\section{\textbf{Achievements / Competitions}}

 \resumeSubHeadingListStart

          \resumeProjectHeading
          {\footnotesize\emph{Algorithmic programming contests, all entered with team \textbf{Five4Five}:}}{}
          \vspace{-6pt}

          \resumeProjectHeading
          {\textbf{\href{https://ieeextreme.org/ieeextreme-17-0-ranking/}{\underline{IEEEXtreme 17.0}}} $|$ \footnotesize{24-hour global --- \textbf{Global Rank 174} of 7,091}}{\textcolor{MidnightBlue}{Nov. 2023}}
          \vspace{-6pt}

          \resumeProjectHeading
          {\textbf{MoraXtream 8.0} $|$ \footnotesize{12-hour, IEEE Student Branch, Univ. of Moratuwa --- \textbf{National Rank 4} of 450+}}{\textcolor{MidnightBlue}{Nov. 2023}}
          \vspace{-6pt}

          \resumeProjectHeading
          {\textbf{Aces Coders v10.0} $|$ \footnotesize{12-hour, ACES, Univ. of Peradeniya --- \textbf{National Rank 8} of 350+}}{\textcolor{MidnightBlue}{Oct. 2023}}
          \vspace{-6pt}

          \resumeProjectHeading
          {\textbf{Aces PreCoders v10.0} $|$ \footnotesize{6-hour --- \textbf{Rank 2} of 50+}}{\textcolor{MidnightBlue}{Sep. 2023}}
          \vspace{-6pt}

          \resumeProjectHeading
          {\textbf{NBQSA} $|$ \footnotesize{National ICT Awards --- IntelliSwitcher: Intelligent Domestic Energy Optimizing System $|$ \href{https://github.com/IntelliSwitcher}{\faIcon{github}}}}{\textcolor{MidnightBlue}{Aug. 2023}}
          \vspace{-3pt}
    \resumeSubHeadingListEnd

% -------------------- Certifications --------------------
\section{\textbf{Selected Certificates}}
    \resumeSubHeadingListStart

          \resumeProjectHeading
          {\textbf{Certificate Course in Molecular Biology \& Biotechnology} $|$ \footnotesize{\textbf{100 hours}}}{\textcolor{MidnightBlue}{2026}}
          \vspace{-14pt}
          \resumeProjectHeading
          {\footnotesize\emph{Agricultural Biotechnology Centre, University of Peradeniya}}{}
          \vspace{-2pt}
          \resumeItemListStart
            \resumeItem{Formal wet-lab and bioinformatics training taken to ground my computational work in experimental biology: \textbf{gene expression regulation and mechanisms}, next-generation and Sanger sequencing, genome organization, recombinant DNA cloning, DNA/RNA extraction, PCR, and gel electrophoresis; bioinformatics modules on NCBI resources, primer design, \textbf{Nanopore transcriptomics analysis}, and molecular phylogenetics. (40h theory, 40h practical, 20h assessed independent work.)}
          \resumeItemListEnd

        \vspace{2pt}
          \resumeProjectHeading
          {\textbf{Machine Learning Specialization} $|$ \footnotesize{Stanford University and DeepLearning.AI(Coursera)}}{\textcolor{MidnightBlue}{Apr. 2024}}
          \resumeItemListStart
            \resumeItem{{\href{https://coursera.org/share/c0d953d1c02645962a825bd697413435}{\underline{Supervised Machine Learning: Regression and Classification}}}}     \resumeItem{{\href{https://www.coursera.org/account/accomplishments/verify/HT4C6JTZZ57R}{\underline{Advanced Learning Algorithms}}}}
        \resumeItemListEnd

        % \vspace{1pt}
        % \resumeProjectHeading
        %   {\textbf{\href{https://courses.edx.org/certificates/3d03dd7fa6814b86b7d65ddc3d1212d4}{AWS Cloud Practitioner Essentials}} $|$ \footnotesize{AWS and edX}}{\textcolor{MidnightBlue}{Oct. 2023}}

        \vspace{1pt}

        \resumeProjectHeading
          {\textbf{\href{https://courses.edx.org/certificates/a2edf86bf72a4bd088f045e94203f3f6}{Python for Data Science, AI and Development}} $|$ \footnotesize{IBM and Coursera}}{\textcolor{MidnightBlue}{Mar. 2023}}



        % \resumeProjectHeading
        %   {\textbf{\href{https://courses.edx.org/certificates/5d016e74d8ce48c1b03aa545b937d5ce}{Linux Commands and Shell Scripting Basics}} $|$ \footnotesize{IBM and edX}}{\textcolor{MidnightBlue}{Feb. 2023}}
    \vspace{4pt}
    \resumeSubHeadingListEnd

% % -------------------- Extra-Curricular Activities --------------------
% \section{\textbf{Extra-Curricular Activitie}s}
% \resumeSubHeadingListStart
%         \vspace{0pt}
%         \resumeProjectHeading
%             { Member of the Web Consultation team of University of Peradeniya\vspace{8pt}}{\textcolor{MidnightBlue}{2023 -- Present}}
%         \vspace{-14pt}
%         \resumeProjectHeading
%             { Senior Member of FASE \small{(Foundation of Astronomical Studies and Exploration)}\vspace{8pt}}{\textcolor{MidnightBlue}{2019 -- Present}}
%         \vspace{-14pt}
%         \resumeProjectHeading
%             { Member of Rotaract Club of University of Peradeniya\vspace{8pt}}{\textcolor{MidnightBlue}{2022 -- Present}}
%         \vspace{-14pt}
%         \resumeProjectHeading
%             { Member of the Athletics team of Dharmaraja College, Kandy\vspace{8pt}}{\textcolor{MidnightBlue}{2013 -- 2017}}
%         \vspace{-4pt}
%     \resumeSubHeadingListEnd


% -------------------- REFERENCES --------------------
\section{\textbf{References}}
\vspace{5pt}
\resumeSubHeadingListStart

\item{
    \small\emph{All three referees are at the University of Peradeniya, Sri Lanka.}
    \\ \vspace{4pt}

    \textbf{\href{https://people.ce.pdn.ac.lk/staff/academic/roshan-ragel/}{\underline{Prof. Roshan G. Ragel}}} $|$ \href{mailto:roshanr@eng.pdn.ac.lk}{roshanr@eng.pdn.ac.lk}
    \\ \vspace{1pt}
    \small{Professor, Department of Computer Engineering}
    \\ \vspace{5pt}

    \textbf{\href{https://people.ce.pdn.ac.lk/staff/academic/damayanthi-herath/}{\underline{Dr. Damayanthi Herath}}}  $|$ \href{mailto:damayanthiherath@eng.pdn.ac.lk}{damayanthiherath@eng.pdn.ac.lk}
    \\ \vspace{1pt}
    \small{Senior Lecturer, Department of Computer Engineering; Coordinator, Computational Biology (CompBio) Group}
    \\ \vspace{1pt}
    \small{Co-author, TARVI (MERCon 2026)}
    \\ \vspace{5pt}

    \textbf{\href{https://roshangodaliyadda888.github.io/}{\underline{Prof. G. M. Roshan I. Godaliyadda}}}  $|$ \href{mailto:roshang@eng.pdn.ac.lk}{roshang@eng.pdn.ac.lk}
    \\ \vspace{1pt}
    \small{Professor, Department of Electrical \& Electronic Engineering}
    \\ \vspace{1pt}
    \small{Deputy Director (Research \& Innovation), MARC}
}

\resumeSubHeadingListEnd

\end{document}
